\section{Stack Direction}

\begin{itemize}
  \item The VIV-32 stack grows downward.
  \item Stack growth proceeds from high addresses toward low addresses.
  \item A push operation decrements \(sp\) before storing data.
  \item A pop operation loads data before incrementing \(sp\).
\end{itemize}

\section{Stack Alignment}

\begin{itemize}
  \item The stack pointer must be 8-byte aligned at every public function call boundary.
  \item The 8-byte alignment rule allows 64-bit ABI values to be placed on the stack with natural alignment.
  \item Future floating-point or vector extensions may define stricter alignment requirements for specific data types.
\end{itemize}

\section{Initial Stack Pointer}

\begin{itemize}
  \item The platform initialises \(sp\) to the first address above the usable stack region.
  \item On the default platform profile, if RAM extends to the byte immediately below the MMIO region, the initial stack pointer may be \texttt{0xFFFF0000}.
\end{itemize}

\section{Frame Pointer}

\begin{itemize}
  \item Register \(fp\), corresponding to \(r13\), may be used as a frame pointer.
  \item Leaf functions and optimised functions may omit the frame pointer.
  \item If a function omits the frame pointer, it may use \(r13\) as a callee-saved general-purpose register.
\end{itemize}

\section{Red Zone}

\begin{itemize}
  \item No red zone is defined.
  \item Memory below \(sp\) is not reserved.
  \item Memory below \(sp\) may be overwritten by interrupts, traps, or called functions.
\end{itemize}

\section{Interrupt and Trap Stack Use}

\begin{itemize}
  \item Interrupt and trap handlers use the current stack in VIV-32 v0.1.
  \item Future privileged execution models may define separate kernel or interrupt stacks.
\end{itemize}

\section{Typical Stack Frame}

\begin{itemize}
  \item A typical non-leaf function stack frame may contain:
  \begin{itemize}
    \item saved link register;
    \item saved frame pointer;
    \item saved callee-saved registers;
    \item local variables;
    \item spilled temporaries;
    \item outgoing stack arguments for nested calls.
  \end{itemize}
  \item The exact layout is determined by the compiler or hand-written assembly routine, provided ABI preservation and alignment rules are satisfied.
\end{itemize}
